\section{Experiments and Analysis}
\subsection{Experimental Setup}
\label{sec:setup}
The evaluation of zero-shot fault diagnosis models requires a rigorous experimental design that accounts for the disparity between semantic descriptions and physical sensor measurements. In this section, we delineate the comprehensive experimental framework employed to evaluate the proposed Diff-LM-GZSL framework. This includes a detailed description of the benchmark datasets, the state-of-the-art baseline models used for comparison, the evaluation metrics for zero-shot (ZSL) and generalized zero-shot (GZSL) scenarios, and the specific implementation details of our novel language-driven diffusion-based generative model. Our setup is designed to test the model's ability to generalize from high-level text-based fault descriptions provided by maintenance manuals or domain experts to the raw, noisy industrial sensor data across multiple physical domains.

\subsubsection{Benchmark Datasets}
To validate the robustness, scalability, and generalizability of the Diff-LM-GZSL framework, we utilize three prominent industrial fault diagnosis datasets: the Tennessee Eastman Process (TEP) \cite{Downs1993TEP}, the Hydraulic System dataset \cite{Helwig2015hydraulic}, and the Case Western Reserve University (CWRU) bearing dataset \cite{Smith2015CWRU}. These datasets represent different failure modes (process-based vs. component-based) and signal characteristics (continuous low-frequency vs. high-frequency vibration signals).

\begin{enumerate}
    \item \textbf{TEP Dataset}: A standard benchmark for process monitoring and control research, the TEP dataset is a simulation of a real-world chemical plant. It comprises 52 variables, including 41 measured variables (recorded at 1.5 to 3-minute intervals) and 11 manipulated variables. It includes 21 types of programmed faults in addition to the normal operating state. The variables represent a mix of continuous and discrete measurements, presenting a high-dimensional, non-linear monitoring challenge that is characteristic of modern chemical production.
    \item \textbf{Hydraulic System Dataset}: This dataset involves a modular multi-sensor rig monitoring 7 distinct fault conditions across various hydraulic components including the cooler, valve, pump, and hydraulic accumulator. It provides a complex multi-rate sensor environment where different sensors (pressure, temperature, flow, vibration) sampled at frequencies ranging from 1Hz to 100Hz must be fused for accurate diagnosis. This heterogeneous sensor setup is a critical test for the alignment capabilities of our diffusion model.
    \item \textbf{CWRU Bearing Dataset}: A widely used vibration signal dataset involving 10 fault types (inner race, outer race, and ball faults) under 4 different motor loads (0 to 3 hp). These high-frequency vibration signals are transformed into time-frequency spectrograms using a Short-Time Fourier Transform (STFT) to preserve spatial-temporal features during the feature extraction process, allowing the model to process sensor data as 2D visual representations.
\end{enumerate}
Detailed statistics regarding the partitioning of seen/unseen categories, the absolute number of samples per class, and the sensor configurations for each case study are summarized in Table~\ref{tab:datasets}.

\begin{table}[htbp]
\centering
\caption{Comparison of traditional binary attributes and the proposed text descriptions for fault characterization.}
\label{tab:attribute_comparison}
\begin{tabular}{lll}
\hline
Fault Type & Binary Attribute & Proposed Text Description \\
\hline
Inner Race & [1,0,1,0,0,1] & Inner race defect with periodic high-frequency impulse and modulation... \\
Outer Race & [0,1,0,1,1,0] & Outer race fault characterized by low-frequency periodic impacts... \\
Ball Fault & [0,0,1,1,0,1] & Ball element damage resulting in non-stationary spectral components... \\
Valve Sticking & [1,1,0,0,1,0] & Mechanical friction in control valve causing delayed process response... \\
\hline
\end{tabular}
\end{table}

\subsubsection{Baseline Methods for Comparison}
We compare the performance of Diff-LM-GZSL against nine representative ZSL and GZSL methods, which have been selected to represent the technological evolution of the field from traditional attribute transfer approaches to modern deep generative models:
\begin{itemize}
    \item \textbf{FDAT (Fault Diagnosis via Attribute Transfer)}: A pioneer method that uses a manually defined binary attribute matrix to transfer knowledge from seen to unseen fault types. It relies heavily on the quality of manual attribute engineering.
    \item \textbf{SCE (Semantic Consistency Embedding)}: This approach enforces structural consistency between the visual feature space and the semantic space during the embedding process to improve the mapping reliability for unseen classes.
    \item \textbf{GLA-ZSL (Global-Local Alignment)}: Focuses on aligning localized parts of the feature space with specific semantic descriptions, reducing the overall domain gap by focusing on salient diagnostic features.
    \item \textbf{FAGAN (Fault Augmentation GAN)}: A specialized GAN architecture designed to handle class-imbalance in industrial datasets by generating synthetic fault features for rare or unseen classes.
    \item \textbf{SRWGAN (Supervised Residual WGAN)}: A Wasserstein GAN variant that incorporates supervised residual learning to stabilize the training process and improve the fidelity of the generated samples.
    \item \textbf{VAEGAN-AR (VAE-GAN with Absolute-Relative Descriptors)}: A hybrid model that combines the manifold learning capabilities of VAEs with the realistic generation of GANs, using relative descriptors to characterize unseen classes.
    \item \textbf{FREE (Feature Refinement for GZSL)}: Utilizes a self-supervised feature refinement module to mitigate the inherent bias towards seen classes during the GZSL inference phase.
    \item \textbf{GZSLCFD (GZSL for Chemical Fault Diagnosis)}: A method specifically optimized for chemical processes, utilizing localized feature-driven descriptions to handle the high dimensionality of process data.
    \item \textbf{DP-CDDPM-AC (Diffusion-based Class-Conditioned Model)}: A recent diffusion-based approach that uses simple class indicators for the denoising process, serving as the primary generative baseline for our diffusion-based architecture.
\end{itemize}

\subsubsection{Evaluation Metrics}
Following standard theoretical protocols in zero-shot learning literature, performance is evaluated using three primary quantitative metrics. For the standard ZSL scenario (where testing is restricted strictly to unseen classes), we report the Top-1 Accuracy ($acc_u$). For the more realistic and challenging GZSL scenario, where the classifier must handle the entire label space (seen + unseen), we report:
\begin{enumerate}
    \item \textbf{$acc_s$ (Seen Accuracy)}: Percent accuracy on test samples belonging to seen classes during GZSL inference.
    \item \textbf{$acc_u$ (Unseen Accuracy)}: Percent accuracy on test samples belonging to unseen classes during GZSL inference.
    \item \textbf{$H$ (Harmonic Mean)}: Calculated as $H = (2 \times acc_s \times acc_u) / (acc_s + acc_u)$. This is the most crucial metric for GZSL as it penalizes models that perform excessively well on seen classes but fail to generalize to unseen ones due to class-bias.
\end{enumerate}
Additionally, we report the Macro-F1 score to provide a comprehensive view of performance across all classes, ensuring that the results are not skewed by potential class imbalances in the test sets.

\subsubsection{Technical Implementation Details}
The feature extraction architecture is tailored to the specific characteristics of the input data: we utilize a 1D-CNN with five convolutional layers for the TEP and Hydraulic signals, and a ResNet-18 architecture for the CWRU spectrograms. For our language-driven semantic module, we leverage the pre-trained \texttt{bert-base-uncased} model. Detailed textual descriptions of each fault—sourced from industrial diagnostic manuals—are tokenized and passed through BERT to obtain 768-dimensional contextual embeddings. These embeddings are then refined via a two-layer MLP to form the 128-dimensional continuous semantic space.

The core Conditional Latent Diffusion Model (CLDM) is configured with $T=1000$ diffusion steps during both training and sampling. We employ a linear variance schedule starting from $\beta_1=10^{-4}$ up to $\beta_T=0.02$. The denoising network $\epsilon_\theta$ is a deep MLP consisting of four hidden layers with 1024 units each, utilizing LeakyReLU activations, batch normalization, and dropout (p=0.2) for regularization. Cross-modal conditioning is implemented via a multi-head cross-attention mechanism that injects the BERT embeddings into the latent states of the diffusion model at each denoising step. The entire multi-objective framework is trained for 500 epochs using the Adam optimizer with a constant learning rate of $1 \times 10^{-4}$ and a batch size of 128. These experiments were conducted on a high-performance workstation equipped with an NVIDIA RTX 4090 GPU and 64GB of RAM to ensure computational efficiency.

\subsection{Case I: Tennessee Eastman Process (TEP)}
\label{sec:case1}

\subsubsection{Comprehensive Dataset Description and Experimental Grouping}
The Tennessee Eastman Process (TEP) is a highly complex and non-linear chemical process simulation of a realistic industrial plant. The process is characterized by strong coupling between multiple variables and significant noise, making it a gold standard for testing fault diagnosis algorithms. The dataset includes 52 variables: 41 measured variables (such as temperatures, pressures, and flow rates, recorded at 3-minute intervals) and 11 manipulated variables (valve positions and set points, recorded at 1.5-minute intervals). The standard TEP benchmark defines 21 distinct fault types, ranging from simple step changes and oscillations in process variables to slow drifts and valve sticking.

In our zero-shot configuration, we partition the 21 faults into a training set of 15 seen classes and a separate test set of 6 unseen classes. The unseen classes are strategically chosen to represent diverse failure mechanisms: Fault 3 (a step change in the feed temperature), Fault 5 (a step change in the condenser cooling water inlet temperature), Fault 9 (random variation in the feed temperature), Fault 15 (a valve sticking condition in the condenser cooling water), Fault 16 (a physically unknown fault defined by the original TEP simulator), and Fault 21 (a valve sticking condition in the feed flow). This diverse selection ensures that the model cannot simply interpolate between seen classes but must rely on the language-driven semantic descriptions to understand the underlying physical characteristics of the unseen faults. We use 800 samples per class for training the feature extractor and the latent diffusion model, while 200 samples per class are strictly reserved for evaluation.

\subsubsection{GZSL Quantitative Results and Performance Analysis}
The comparative performance for the TEP dataset is summarized in Table~\ref{tab:tep_results}. Our proposed Diff-LM-GZSL framework demonstrates a substantial and statistically significant performance gain over all existing methods. In the generalized zero-shot learning (GZSL) scenario, our model achieves a Harmonic Mean ($H$) of 89.6\%, which represents a 4.2\% absolute improvement over the best-performing generative baseline, DP-CDDPM-AC. Even more importantly, the accuracy on unseen classes ($acc_u$) stands at an impressive 85.2\%, compared to only 80.2\% for the nearest competitor.

The significant margin is primarily attributable to the use of continuous semantic embeddings instead of binary attributes. Traditional methods (e.g., FDAT or SCE) often fail on the TEP dataset because binary attributes are too coarse to distinguish between faults that occur in similar physical locations but have different dynamics (e.g., distinguishing a "step change" from "random variation" in the same reactor). By using BERT-encoded language descriptions, Diff-LM-GZSL captures the fine-grained semantic nuance that a "valve sticking" fault has a distinct temporal and magnitude characteristic from other physical disruptions. The diffusion model then translates this high-dimensional semantic knowledge into a high-fidelity distribution of synthetic features that perfectly match the manifold of real faults.

\begin{table}[htbp]
\centering
\caption{Quantitative Performance and Accuracy Comparison on the TEP Dataset (Case I)}
\label{tab:tep_results}
\begin{tabular}{lcccc}
\hline
Method & $acc_s$ (\%) & $acc_u$ (\%) & $H$ (\%) & Macro-F1 \\
\hline
FDAT & $72.4 \pm 1.2$ & $48.5 \pm 2.1$ & $58.1$ & $0.59$ \\
SCE & $75.1 \pm 0.8$ & $52.3 \pm 1.5$ & $61.6$ & $0.62$ \\
GLA-ZSL & $78.9 \pm 0.9$ & $56.8 \pm 1.8$ & $66.0$ & $0.67$ \\
FAGAN & $81.2 \pm 1.1$ & $65.4 \pm 2.2$ & $72.5$ & $0.73$ \\
SRWGAN & $83.5 \pm 0.7$ & $69.2 \pm 1.9$ & $75.7$ & $0.76$ \\
VAEGAN-AR & $85.0 \pm 1.0$ & $72.1 \pm 1.4$ & $78.0$ & $0.79$ \\
FREE & $88.2 \pm 0.6$ & $79.5 \pm 1.2$ & $83.6$ & $0.84$ \\
GZSLCFD & $90.5 \pm 0.5$ & $78.2 \pm 1.6$ & $83.9$ & $0.85$ \\
DP-CDDPM-AC & $91.3 \pm 0.4$ & $80.2 \pm 1.1$ & $85.4$ & $0.87$ \\
\textbf{Diff-LM-GZSL} & $\mathbf{94.6 \pm 0.3}$ & $\mathbf{85.2 \pm 0.8}$ & $\mathbf{89.6}$ & $\mathbf{0.91}$ \\
\hline
\end{tabular}
\end{table}

\subsubsection{Multi-modal Visualization and Feature Distribution Evaluation}
To gain deeper insights into the diagnostic behavior, we examine the confusion matrix for the TEP unseen classes in Fig.~5. The matrix indicates that our model maintains high classification precision for nearly all unseen faults, with very few misclassifications. Only a minor confusion exists between Fault 15 and Fault 16. This is physically understandable, as Fault 16 is an "unknown" disturbance in the original TEP simulator that happens to exhibit some overlapping thermal characteristics with the condenser cooling water system. 

Furthermore, we utilize t-SNE visualization (Fig.~8) to project the high-dimensional latent space of the model into two dimensions. The visualization reveals the distribution of real seen and unseen features alongside the synthetic samples generated by our conditional diffusion model. The synthetic features for unseen classes (represented in distinct colors) perfectly overlap with the real latent distributions, forming compact and well-separated clusters. This robust alignment confirms that our semantic module and alignment loss successfully mitigate the domain shift problem, ensuring that the synthesized training data is physically consistent with real world sensor data for the classifier to learn from.

\subsection{Case II: Hydraulic System}
\label{sec:case2}

\subsubsection{Detailed Dataset Description and Heterogeneous Sensor Grouping}
The Hydraulic System dataset is collected from a modular industrial test rig where specific component-level faults can be induced under controlled conditions. The rig consists of a primary circuit powered by a motor and pump, and a secondary circuit used for load simulation. Continuous sensors monitor a wide array of physical quantities such as pressure (at 6 locations), flow rates (2 locations), temperature (4 locations), motor power, and vibration levels. The data is characterized by multi-rate sampling—a common challenge in industrial monitoring—requiring an advanced preprocessing stage to align the 1Hz temperature data with high-frequency 100Hz pressure signals using interpolation and normalization techniques.

In our zero-shot experimental design, we categorize the rig's condition into 7 distinct fault states: Cooler condition (3 states), Valve condition (4 states), Internal leakage (3 states), and Accumulator status (4 states). Following a rigorous 5-seen/2-unseen class split, the unseen conditions are selected as the "Reduced Cooler Efficiency" and "Internal Pump Leakage (Severe)" classes. These specific classes were chosen because their physical symptoms are often subtle and can be masked by the system's internal feedback control loops, making them particularly difficult to diagnose without existing training data. The primary objective is to determine if the language-driven module can correctly translate a textual description of "reduced thermal efficiency" into the specific multi-modal patterns captured by the heterogenous sensors.

\subsubsection{Experimental Results and Cross-System Consistency Analysis}
The performance of Diff-LM-GZSL on the Hydraulic dataset is presented in Table~\ref{tab:hydraulic_results}. We observe a consistent trend that reinforces the findings from the TEP case study, where our method outperforms the state-of-the-art baselines by a wide margin. The Harmonic Mean $H$ reaches an impressive 92.4\%, and the Macro-F1 score stands at 0.93. This represents a 4.9\% improvement in $H$ compared to the nearest diffusion-based baseline, demonstrating that our method is exceptionally robust for diagnostic systems with heterogeneous and multi-rate sensor configurations.

The high accuracy on unseen classes (89.7\%) is particularly noteworthy for a hydraulic system. It suggests that the BERT-based semantic encoder is capable of successfully modeling the deep physical relationships—for instance, the causal connection between thermal "leakage" and specific pressure drop patterns—allowing the latent diffusion model to generate accurate feature representations even for severe failure modes that have never been physically observed during the model's training phase. The confusion matrix for the hydraulic system is provided in Fig.~\ref{fig:cm_hydraulic}, highlighting the precise boundaries between fault states.

\begin{table}[htbp]
\centering
\caption{Quantitative Performance on the Hydraulic System Dataset (Case II)}
\label{tab:hydraulic_results}
\begin{tabular}{lcccc}
\hline
Method & $acc_s$ (\%) & $acc_u$ (\%) & $H$ (\%) & Macro-F1 \\
\hline
FDAT & $75.2 \pm 1.4$ & $50.1 \pm 2.5$ & $60.1$ & $0.61$ \\
SCE & $76.8 \pm 1.1$ & $55.4 \pm 2.0$ & $64.4$ & $0.65$ \\
GLA-ZSL & $80.5 \pm 1.0$ & $60.2 \pm 2.2$ & $68.9$ & $0.70$ \\
FAGAN & $83.4 \pm 1.5$ & $70.5 \pm 2.8$ & $76.4$ & $0.78$ \\
SRWGAN & $85.6 \pm 0.9$ & $74.3 \pm 2.1$ & $79.5$ & $0.81$ \\
VAEGAN-AR & $86.2 \pm 1.2$ & $76.8 \pm 1.5$ & $81.2$ & $0.83$ \\
FREE & $89.4 \pm 0.8$ & $82.4 \pm 1.4$ & $85.8$ & $0.87$ \\
GZSLCFD & $91.2 \pm 0.6$ & $81.5 \pm 1.7$ & $86.1$ & $0.88$ \\
DP-CDDPM-AC & $92.5 \pm 0.5$ & $83.1 \pm 1.2$ & $87.5$ & $0.89$ \\
\textbf{Diff-LM-GZSL} & $\mathbf{95.3 \pm 0.4}$ & $\mathbf{89.7 \pm 0.7}$ & $\mathbf{92.4}$ & $\mathbf{0.93}$ \\
\hline
\end{tabular}
\end{table}

\subsubsection{Extensive Ablation Study on Semantic Types and Core Components}
To further validate the individual design choices of the Diff-LM-GZSL framework, we conducted a series of rigorous ablation experiments on the Hydraulic System dataset. These tests, summarized in Table~\ref{tab:ablation}, highlight the contribution of three fundamental modules: Large Language Model (LLM) continuous semantics, the Latent Diffusion generative process, and the Semantic Alignment module.

\begin{table}[htbp]
\centering
\caption{Ablation Study Results on the Hydraulic System Dataset (Case II)}
\label{tab:ablation}
\begin{tabular}{lcccc}
\hline
Configuration & $acc_s$ (\%) & $acc_u$ (\%) & $H$ (\%) \\
\hline
w/o LLM Semantics (Binary Attributes Only) & $91.4$ & $74.2$ & $81.9$ \\
w/o Diffusion (Using GAN Generator) & $86.5$ & $78.1$ & $82.1$ \\
w/o Semantic Alignment (No MMD Loss) & $92.8$ & $84.5$ & $88.5$ \\
\textbf{Full Diff-LM-GZSL} & $\mathbf{95.3}$ & $\mathbf{89.7}$ & $\mathbf{92.4}$ \\
\hline
\end{tabular}
\end{table}

As the results indicate, the most significant performance degradation occurs when removing the LLM semantics. Substituting our continuous BERT embeddings with traditional binary attributes causes the unseen class accuracy to plummet from 89.7\% to 74.2\%, a reduction of over 15\%. This clearly illustrates that binary attributes do not provide sufficient diagnostic information for zero-shot generalization in complex hydraulic systems. Furthermore, replacing the diffusion process with a standard GAN generator results in significantly lower accuracy across both seen and unseen classes, likely due to the superior high-dimensional manifold-modeling capability of diffusion models. Finally, removing the semantic alignment loss leads to a 3.9\% decrease in the Harmonic Mean, confirming its essential role in reducing the visual-semantic domain gap. These findings underscore that the synergy between language models and latent diffusion probabilistic models is the cornerstone for reliable zero-shot fault diagnosis.
